% Style follows ../class2/class2PT26.tex.
% Build in this directory (run twice to resolve page references and bookmarks):
% pdflatex -interaction=nonstopmode -halt-on-error class4PT26.tex
% Lecture topics follow PT26w/lectureNote.tex, sections 4 and 5.
\documentclass[a4paper,10pt,pdflatex,ja=standard,jafont=ipaex-type1,
  magstyle=real,scale=1]{bxjsarticle}
\usepackage{lmodern}
\usepackage{amsmath,amssymb,mathtools}
\usepackage{geometry}
\geometry{reset,left=22mm,right=22mm,top=20mm,bottom=17mm,
  headheight=13pt,headsep=6mm,footskip=10mm}
\usepackage{xcolor}
\usepackage{enumitem}
\usepackage{array,booktabs,tabularx}
\usepackage{fancyhdr}
\usepackage{lastpage}
\usepackage{hyperref}

\definecolor{ink}{HTML}{193A5C}
\definecolor{accent}{HTML}{237D82}
\definecolor{soft}{HTML}{EAF3F3}
\definecolor{muted}{HTML}{52616B}
\newcommand{\coursename}{確率論と統計学 I}
\newcommand{\lecturenumber}{4}
\hypersetup{unicode=true,colorlinks=true,urlcolor=accent,linkcolor=ink,
  pdfauthor={中島秀太},pdftitle={第4・5回：独立性と共分散、代表的な確率分布}}
\pdfstringdefDisableCommands{\def\quad{ }}
\pagestyle{fancy}
\fancyhf{}
\fancyhead[L]{\small\color{ink}\coursename{}・授業要約と演習}
\fancyhead[R]{\small\color{ink}第\lecturenumber 回}
\fancyfoot[C]{\small\color{ink}\thepage/\pageref*{LastPage}}
\renewcommand{\headrulewidth}{0.3pt}
\setlength{\parindent}{0pt}
\setlength{\parskip}{2pt}
\setlength{\abovedisplayskip}{4pt}
\setlength{\belowdisplayskip}{4pt}
\setlength{\abovedisplayshortskip}{3pt}
\setlength{\belowdisplayshortskip}{4pt}
\setlist[itemize]{leftmargin=1.8em,itemsep=1pt,topsep=2pt,parsep=0pt}
\setlist[enumerate]{leftmargin=2em,itemsep=3pt,topsep=3pt,parsep=0pt}
\newcommand{\topic}[1]{\par\vspace{5pt}
  {\large\sffamily\mdseries\color{accent}#1}\par\vspace{2pt}}
\newcommand{\handouttitle}[1]{%
  \pdfbookmark[0]{#1}{lecture\lecturenumber}%
  \begin{center}
    {\Large\color{ink}#1}\par\vspace{3mm}
    {\color{accent}\rule{0.88\linewidth}{0.7pt}}
  \end{center}\vspace{-2mm}}
\newcommand{\keybox}[1]{\par\vspace{3pt}\begingroup
  \setlength{\fboxsep}{5pt}
  \colorbox{soft}{\parbox{\dimexpr\linewidth-2\fboxsep\relax}{#1}}%
  \endgroup\par\vspace{2pt}}
\newcommand{\examplelabel}[1]{{\sffamily\mdseries\color{ink}#1}}
\newcommand{\F}{\mathcal F}
\newcommand{\R}{\mathbb R}
\newcommand{\E}{\mathbb E}
\newcommand{\Var}{\operatorname{Var}}
\newcommand{\Cov}{\operatorname{Cov}}

\begin{document}
\fontsize{10pt}{13.5pt}\selectfont
\handouttitle{第4回\quad 独立性と共分散}

\topic{確率変数の独立性}
$(\Omega,\F,P)$ を確率空間、$\mathcal B(\R)$ を $\R$ のボレル集合族とする。
確率変数 $X,Y:\Omega\to\R$ が独立であるとは、
すべての $A,B\in\mathcal B(\R)$ に対して次が成り立つことである。
$$
  P(X\in A,\,Y\in B)=P(X\in A)P(Y\in B).
$$
ここで左辺は $P(\{X\in A\}\cap\{Y\in B\})$ を表す。
離散型の場合は、すべての取り得る値 $x,y$ について
$$
  P(X=x,\,Y=y)=P(X=x)P(Y=y)
$$
が成り立つことと同値である。

\topic{独立性の保存と積の期待値}
$X,Y$ が独立で $f,g:\R\to\R$ がボレル可測なら、$f(X),g(Y)$ も独立である。
さらに $\E[|f(X)|],\E[|g(Y)|]<\infty$ なら、その積も可積分で
$$
  \E[f(X)g(Y)]=\E[f(X)]\E[g(Y)].
$$
特に、独立な可積分確率変数について $\E[XY]=\E[X]\E[Y]$ となる。

期待値の線形性 $\E[X+Y]=\E[X]+\E[Y]$ には独立性は不要である。

\topic{分散と共分散}
$\E[X^2],\E[Y^2]<\infty$ を仮定する。分散と共分散を次で定義する。
$$
  \begin{aligned}
    \Var(X)&:=\E[(X-\E[X])^2]=\E[X^2]-(\E[X])^2,\\
    \Cov(X,Y)&:=\E[(X-\E[X])(Y-\E[Y])]=\E[XY]-\E[X]\E[Y].
  \end{aligned}
$$
$a,b\in\R$ に対して
$$
  \Var(aX+bY)=a^2\Var(X)+2ab\Cov(X,Y)+b^2\Var(Y).
$$
\keybox{\examplelabel{要点：}上の二次モーメントの仮定のもとで、独立なら無相関
（$\Cov(X,Y)=0$）である。逆は一般に成り立たない。
$\Var(X+Y)=\Var(X)+\Var(Y)$ は無相関であれば成り立つ。}

\topic{対独立と相互独立}
確率変数 $X_1,\ldots,X_n$（$n\ge2$ は整数）が相互独立であるとは、
次がすべての $B_1,\ldots,B_n\in\mathcal B(\R)$ に対して成り立つことである。
$$
  P\left(\bigcap_{i=1}^n\{X_i\in B_i\}\right)
  =\prod_{i=1}^nP(X_i\in B_i)
$$
異なる添字のどの2変数も独立なとき、対独立という。
相互独立なら（一部の $B_i$ を $\R$ として）対独立が従うが、逆は一般に成り立たない。

\topic{演習問題}
\begin{enumerate}[label=\arabic*.]
  \item 公平なサイコロを独立に2回投げ、その出目を $X,Y$ とする。
  $\E[X+Y]$、$\Var(X+Y)$、$\Cov(X+Y,X-Y)$ を求めよ。
  \item $\Omega=[-1,1]$、$\F=\mathcal B([-1,1])$、$P=\lambda/2$（$\lambda$ はルベーグ測度）とし、
  $X(\omega)=\omega$、$Y(\omega)=\omega^2$ とする。$\Cov(X,Y)$ を求めよ。
  事象 $\{|X|\le1/2\}$ と $\{Y\le1/4\}$ を用い、$X,Y$ が独立か判定せよ。
\end{enumerate}

\topic{レポート問題}
\begin{enumerate}[label=\arabic*.]
  \item 独立な $X,Y$ がそれぞれ $-1,1$ を確率 $1/2$ でとるとし、$Z=XY$ とおく。
  各ペアの同時分布を求め、$X,Y,Z$ が対独立であることを示せ。
  次に $\{X=1,Y=1,Z=1\}$ の確率を用い、相互独立ではないことを示せ。
  \item $X,Y$ を独立な確率変数、$f,g:\R\to\R$ をボレル可測関数とする。
  集合の逆像を用いて、$f(X),g(Y)$ が独立であることを示せ。
  \item これまでの授業や演習の感想または改善してほしいところなど。
\end{enumerate}

\newpage
\renewcommand{\lecturenumber}{5}
\handouttitle{第5回\quad 代表的な確率分布}

\topic{分布から期待値を計算する}
$\mu_X(B):=P(X\in B)$（$B\in\mathcal B(\R)$）を $X$ の分布とする。
ボレル可測関数 $g:\R\to\R$ について、
$g\ge0$ または $\E[|g(X)|]<\infty$ なら
$$
  \E[g(X)]=\int_\R g(x)\,\mu_X(dx).
$$
離散分布 $P(X=x_k)=q_k$（$x_k$ は相異なり、$q_k\ge0$、$\sum_k q_k=1$）なら
$$
  \E[g(X)]=\sum_k g(x_k)q_k\text{。}
$$
非負ボレル関数 $f$ が $\mu_X(B)=\int_B f(x)\,dx$ を満たすとき、$f$ を確率密度といい、
$\E[g(X)]=\int_\R g(x)f(x)\,dx$ となる。いずれも上と同じ条件で用いる。
特に $g(x)=x,x^2$ とすれば、期待値と二次モーメントを計算できる。

\topic{離散分布：成功回数・待ち時間・発生回数}
$n\ge1$ は整数、$0<p<1$、$\lambda>0$ とする。

\begingroup
\setlength{\tabcolsep}{5pt}
\renewcommand{\arraystretch}{1.55}
\begin{tabularx}{\linewidth}{@{}l c >{\centering\arraybackslash}X c c@{}}
\toprule
\examplelabel{分布} & \examplelabel{$k$ の範囲} & \leavevmode\examplelabel{$P(X=k)$}
  & \examplelabel{$\E[X]$} & \examplelabel{$\Var(X)$} \\
\midrule
ベルヌーイ $\mathrm{Ber}(p)$ & $0,1$ & $p^k(1-p)^{1-k}$ & $p$ & $p(1-p)$ \\
二項 $\mathrm{Bin}(n,p)$ & $0,\ldots,n$ & $\binom nk p^k(1-p)^{n-k}$ & $np$ & $np(1-p)$ \\
幾何 $\mathrm{Geo}(p)$ & $1,2,\ldots$ & $(1-p)^{k-1}p$ & $\dfrac1p$ & $\dfrac{1-p}{p^2}$ \\
ポアソン $\mathrm{Poi}(\lambda)$ & $0,1,\ldots$ & $e^{-\lambda}\dfrac{\lambda^k}{k!}$ & $\lambda$ & $\lambda$ \\
\bottomrule
\end{tabularx}
\par\endgroup
\smallskip
ベルヌーイ分布は1回の試行の成功を1、失敗を0で表す。
$X_1,\ldots,X_n\sim\mathrm{Ber}(p)$ が独立なら、その和は $\mathrm{Bin}(n,p)$ に従う。

幾何分布は初めて成功するまでの試行回数（成功した回を含む）を表す。
ポアソン分布は、一定の区間内の発生回数のモデルとして用いる。

\topic{連続分布：一様・指数分布}
$a<b$、$\lambda>0$ とする。
$1_A(x)$ は $x\in A$ のとき1、それ以外で0をとる指示関数である。

\begingroup
\setlength{\tabcolsep}{5pt}
\renewcommand{\arraystretch}{1.75}
\begin{tabularx}{\linewidth}{@{}l >{\centering\arraybackslash}X c c@{}}
\toprule
\examplelabel{分布} & \leavevmode\examplelabel{密度 $f(x)$（$x\in\R$）}
  & \examplelabel{$\E[X]$} & \examplelabel{$\Var(X)$} \\
\midrule
一様 $U(a,b)$ & $\dfrac{1}{b-a}\,1_{[a,b]}(x)$ & $\dfrac{a+b}{2}$ & $\dfrac{(b-a)^2}{12}$ \\
指数 $\mathrm{Exp}(\lambda)$ & $\lambda e^{-\lambda x}\,1_{[0,\infty)}(x)$ & $\dfrac1\lambda$ & $\dfrac1{\lambda^2}$ \\
\bottomrule
\end{tabularx}
\par\endgroup
\smallskip
指数分布の $\lambda$ は率のパラメータであり、平均待ち時間は $1/\lambda$ である。

\topic{演習問題}
\begin{enumerate}[label=\arabic*.]
  \item $S\sim\mathrm{Bin}(5,1/3)$ とする。
  $P(S=0)$、$P(S=2)$、$\E[S]$、$\Var(S)$ を求めよ。
  \item 成功確率 $p=1/4$ の試行を独立に繰り返し、初成功までの試行回数を $T$ とする。
  $P(T>3)$、$\E[T]$、$\Var(T)$ を求めよ。
  \item $W\sim\mathrm{Exp}(2)$ とする。密度を積分して $P(W>1)$ を求めよ。
  また、$\E[W]$、$\Var(W)$ を求めよ。
\end{enumerate}

\topic{レポート問題}
\begin{enumerate}[label=\arabic*.]
  \item $\lambda,\nu>0$ とし、独立な $X\sim\mathrm{Poi}(\lambda)$、$Y\sim\mathrm{Poi}(\nu)$ に対して
  $X+Y\sim\mathrm{Poi}(\lambda+\nu)$ を示せ。
  $P(X+Y=k)=\sum_{j=0}^kP(X=j,\,Y=k-j)$ を用いてよい。
\end{enumerate}
\end{document}
